\documentclass[12pt, a4paper]{article}

% --- Pakiety ---
\usepackage[utf8]{inputenc}
\usepackage[T1]{fontenc}
\usepackage[polish]{babel}
\usepackage{amsmath, amssymb}
\usepackage{graphicx}
\usepackage{float}
\usepackage{geometry}
\usepackage{hyperref}
\usepackage{booktabs}
\usepackage{caption}
\usepackage{listings}
\usepackage{xcolor}
\usepackage{parskip}

\geometry{
  a4paper,
  left=2.5cm,
  right=2.5cm,
  top=2.5cm,
  bottom=2.5cm
}

% --- Zakaz dzielenia wyrazów ---
\hyphenpenalty=10000
\exhyphenpenalty=10000
\emergencystretch=3em

\hypersetup{
  colorlinks=true,
  linkcolor=blue,
  urlcolor=blue,
  citecolor=blue
}

% --- Zmiana nazwy float figure: Rysunek -> Wykres ---
\addto\captionspolish{\renewcommand{\figurename}{Wykres}}

% --- Styl listingów ---
\lstset{
  language=Python,
  basicstyle=\ttfamily\small,
  keywordstyle=\color{blue}\bfseries,
  commentstyle=\color{gray}\itshape,
  stringstyle=\color{orange!80!black},
  numbers=left,
  numberstyle=\tiny\color{gray},
  numbersep=8pt,
  frame=single,
  breaklines=true,
  breakatwhitespace=true,
  captionpos=b,
  showstringspaces=false,
  columns=flexible,
  keepspaces=true,
  inputencoding=utf8,
  extendedchars=true,
  literate={ą}{{\k{a}}}1 {ę}{{\k{e}}}1 {ó}{{\'{o}}}1
            {ś}{{\'{s}}}1 {ł}{{\l{}}}1 {ź}{{\'{z}}}1
            {ż}{{\.z}}1 {ć}{{\'{c}}}1 {ń}{{\'{n}}}1
            {Ą}{{\k{A}}}1 {Ę}{{\k{E}}}1 {Ó}{{\'{O}}}1
            {Ś}{{\'{S}}}1 {Ł}{{L}}1,
}

% --- Strona tytułowa ---
\title{
  \textbf{Metody Obliczeniowe w Nauce i Technice} \\[0.5em]
  \large Laboratorium 5 \textemdash{} Aproksymacja
}
\author{Mateusz Stoch}
\date{14 kwietnia 2026}

\begin{document}

\maketitle

% ============================================================
\section{Zadanie 1 \textemdash{} Aproksymacja populacji USA}
% ============================================================

% --- 1.1 Opis zadania ---
\subsection{Opis zadania}

Dane wejściowe to populacja Stanów Zjednoczonych w~latach 1900--1980
(9~punktów co 10 lat). Należy:
\begin{enumerate}
  \item[(a)] Wykonać aproksymację średniokwadratową wielomianami stopnia
             $m = 0, 1, \ldots, 6$ i ekstrapolować każdy wielomian do
             roku~1990. Porównać wyniki z~prawdziwą wartością
             248\,709\,873 i znaleźć stopień $m$ dający najmniejszy
             błąd względny.
  \item[(b)] Wyznaczyć optymalny stopień~$m$ za pomocą kryterium
             informacyjnego Akaikego ze składnikiem korygującym
             (AICc) i sprawdzić, czy pokrywa się on z~wynikiem
             z~punktu (a).
\end{enumerate}

% --- 1.2 Argumentacja rozwiązania ---
\subsection{Argumentacja rozwiązania}

\subsubsection{Normalizacja argumentu}

Dane obejmują lata 1900--1990. Używanie surowych lat jako argumentów
wielomianu prowadziłoby do ogromnych wartości liczbowych i złego
uwarunkowania macierzy układu. Dlatego argumenty normalizuje się do
przedziału $[-1, 1]$:
\[
  t = \frac{\text{rok} - 1940}{40}.
\]
Rok 1990 odpowiada $t_{1990} = (1990 - 1940) / 40 = 1{,}25$.

\subsubsection{Aproksymacja średniokwadratowa}

Aproksymacja różni się od interpolacji, ponieważ nie wymaga, by wielomian
dokładnie przechodził przez węzły. Szukamy wielomianu $\hat{y}(t)$,
który \emph{minimalizuje} sumę kwadratów odchyleń:
\[
  \min_{\mathbf{c}} \sum_{i=1}^{n} \bigl[y_i - \hat{y}(t_i)\bigr]^2.
\]
W~praktyce sprowadza się to do rozwiązania układu \emph{normalnych
równań} lub rozwiązania zadania metodą najmniejszych kwadratów
(np.~\texttt{np.linalg.lstsq}).

\begin{lstlisting}[caption={Dopasowanie wielomianu stopnia m metodą najmniejszych kwadratów}]
def fit_polynomial_ls(t_data, y_data, m):
    A = np.vander(t_data, N=m + 1, increasing=True)
    c, residuals, rank, sv = np.linalg.lstsq(A, y_data, rcond=None)
    return c
\end{lstlisting}

\subsubsection{Kryterium AICc}

Kryterium informacyjne Akaikego pozwala wybrać kompromis między
dokładnością modelu a jego złożonością. Dla próbki o małym rozmiarze
($n/k < 40$, gdzie $k = m + 1$) stosuje się wersję skorygowaną:
\[
  \text{AIC} = 2k + n \ln\!\left(\frac{\mathrm{SSE}}{n}\right),
  \qquad
  \text{AICc} = \text{AIC} + \frac{2k(k+1)}{n - k - 1},
\]
gdzie $\mathrm{SSE} = \sum_{i=1}^{n}[y_i - \hat{y}(t_i)]^2$.
Mniejsza wartość AICc odpowiada lepszemu modelowi.

% --- 1.3 Wyniki ---
\subsection{Wykresy i wyniki}

\subsubsection{(a) Ekstrapolacja do 1990}

\begin{figure}[H]
  \centering
  \includegraphics[width=0.9\textwidth]{wykres1.png}
  \caption{Wielomiany aproksymujące populację USA (stopnie $m = 0 \ldots 6$)
           z ekstrapolacją do roku~1990. Punkty danych: 1900--1980.}
  \label{fig:aprox_usa}
\end{figure}

\begin{table}[H]
  \centering
  \caption{Ekstrapolacja do roku 1990 -- porównanie wyników.}
  \begin{tabular}{rrr}
    \toprule
    $m$ & Predykcja 1990 & Błąd względny [\%] \\
    \midrule
    0 & 143\,369\,177 & $\approx 42{,}35\%$ \\
    1 & 235\,808\,109 & $\approx 5{,}19\%$  \\
    2 & 254\,712\,945 & $\approx 2{,}41\%$  \\
    3 & 261\,439\,380 & $\approx 5{,}12\%$  \\
    4 & 243\,106\,971 & $\approx 2{,}25\%$  \\
    5 & 220\,442\,802 & $\approx 11{,}37\%$ \\
    6 & 255\,044\,185 & $\approx 2{,}55\%$  \\
    \bottomrule
  \end{tabular}
  \label{tab:ekstrapolacja_usa}
\end{table}

\noindent
Prawdziwa populacja w 1990: 248\,709\,873.\\
\textbf{Najmniejszy błąd względny dla $m = 4$: $\approx 2{,}25\%$.}

\subsubsection{(b) Kryterium AICc}

\begin{figure}[H]
  \centering
  \includegraphics[width=0.9\textwidth]{wykres2.png}
  \caption{Wartości AICc dla wielomianów stopnia $m = 0 \ldots 6$.}
  \label{fig:aicc}
\end{figure}

\begin{table}[H]
  \centering
  \caption{Zestawienie wartości SSE, AIC i~AICc dla każdego stopnia.}
  \begin{tabular}{rrrrrr}
    \toprule
    $m$ & $k{=}m{+}1$ & $n/k$ & SSE & AIC & AICc \\
    \midrule
    0 & 1 & 9{,}00 & $2{,}09 \times 10^{16}$ & 320{,}44 & 321{,}01 \\
    1 & 2 & 4{,}50 & $4{,}10 \times 10^{14}$ & 287{,}06 & 289{,}06 \\
    2 & 3 & 3{,}00 & $8{,}28 \times 10^{13}$ & 274{,}65 & 279{,}45 \\
    3 & 4 & 2{,}25 & $6{,}80 \times 10^{13}$ & 274{,}88 & 284{,}88 \\
    4 & 5 & 1{,}80 & $3{,}51 \times 10^{13}$ & 270{,}93 & 290{,}93 \\
    5 & 6 & 1{,}50 & $2{,}33 \times 10^{13}$ & 269{,}26 & 311{,}26 \\
    6 & 7 & 1{,}29 & $1{,}87 \times 10^{13}$ & 269{,}27 & 381{,}27 \\
    \bottomrule
  \end{tabular}
  \label{tab:aicc_usa}
\end{table}

\noindent
\textbf{Najmniejsze AICc dla $m = 2$ (AICc = 279{,}45).}\\
Najlepszy błąd ekstrapolacji: $m = 4$, wyniki się \textbf{nie pokrywają}.

% --- 1.4 Wnioski ---
\subsection{Wnioski}

\begin{enumerate}
  \item Wielomian stopnia~0 (stała) dał błąd ponad $42\%$, ponieważ
        nie jest w stanie opisać żadnego trendu.

  \item Wielomian liniowy ($m = 1$) znacznie lepiej opisuje dane
        (błąd $\approx 5\%$), ale nadal zbyt prosto odwzorowuje
        krzywoliniowy trend populacji.

  \item Najlepszą predykcję dla roku~1990 dał wielomian stopnia~4
        ($\approx 2{,}25\%$ błędu). Jednocześnie AICc wybiera $m = 2$
        jako kompromis między jakością dopasowania a ryzykiem
        przeuczenia. Rozbieżność wynika z tego, że AICc ``karze''
        modele wieloparametrowe przy małych zbiorach danych, czyli gdy
        $n = 9$ i $k$ rośnie, $n/k$ gwałtownie maleje, a składnik
        karny rośnie.

  \item Modele o wyższych stopniach ($m = 5, 6$) dobrze dopasowują
        dane treningowe (SSE maleje), ale gorzej ekstrapolują (błąd
        rośnie), co jest klasycznym objawem przeuczenia.

  \item AICc stara się ten problem wykryć \emph{bez} zestawu
        testowego, za pomocą kryterium opartego na danych
        treningowych. W~tym zadaniu sprawdza się, ponieważ chroni
        przed wyborem zbyt złożonego modelu.
\end{enumerate}

% ============================================================
\section{Zadanie 2 \textemdash{} Aproksymacja $\sqrt{x}$ wielomianami Czebyszewa}
% ============================================================

% --- 2.1 Opis zadania ---
\subsection{Opis zadania}

Należy wykonać aproksymację ciągłą funkcji $f(x) = \sqrt{x}$
na~przedziale~$[0, 2]$ wielomianem drugiego stopnia, korzystając
z~wielomianów Czebyszewa jako bazy. Cel jest znalezienie wielomianu
\[
  p_2(x) = c_0 T_0(x') + c_1 T_1(x') + c_2 T_2(x'),
\]
gdzie $x' = x - 1$ przelicza $[0,2]$ na~$[-1,1]$,
a~$T_k$ to $k$-ty wielomian Czebyszewa.

% --- 2.2 Argumentacja rozwiązania ---
\subsection{Argumentacja rozwiązania}

\subsubsection{Wielomiany Czebyszewa}

Wielomiany Czebyszewa $T_k(\xi)$ są zdefiniowane rekurencyjnie:
\[
  T_0(\xi) = 1,\quad T_1(\xi) = \xi,\quad
  T_{k+1}(\xi) = 2\xi T_k(\xi) - T_{k-1}(\xi).
\]
Tworzą one bazę ortogonalną na~$[-1,1]$ z wagą
$w(\xi) = 1/\sqrt{1 - \xi^2}$:
\[
  \int_{-1}^{1} T_i(\xi) T_j(\xi)\, \frac{d\xi}{\sqrt{1-\xi^2}}
  = \begin{cases} 0 & i \neq j \\ \pi/2 & i = j \neq 0 \\ \pi & i = j = 0. \end{cases}
\]

\subsubsection{Wyznaczanie współczynników}

Z ortogonalności wynika wzór na współczynniki:
\[
  c_k = \frac{2}{\pi} \int_{-1}^{1} f(x') T_k(\xi)\,
        \frac{d\xi}{\sqrt{1-\xi^2}}, \qquad k > 0,
\]
(dla $c_0$ czynnik $2/\pi$ zastępuje się $1/\pi$). Całkę oblicza
się numerycznie kwadraturą Gaussa--Czebyszewa:
\[
  \int_{-1}^{1} g(\xi) \frac{d\xi}{\sqrt{1-\xi^2}}
  \approx \frac{\pi}{N} \sum_{j=1}^{N} g\!\left(\cos\!\frac{(2j-1)\pi}{2N}\right).
\]

\noindent
Poniżej kluczowe fragmenty implementacji:

\begin{lstlisting}[caption={Rekurencyjne obliczanie wielomianów Czebyszewa}]
def chebyshev_T(k, t):
    if k == 0:
        return np.ones_like(t, dtype=float)
    elif k == 1:
        return np.asarray(t, dtype=float).copy()
    else:
        T_prev2 = np.ones_like(t, dtype=float)
        T_prev1 = np.asarray(t, dtype=float).copy()
        for _ in range(2, k + 1):
            T_curr = 2 * t * T_prev1 - T_prev2
            T_prev2 = T_prev1
            T_prev1 = T_curr
        return T_curr
\end{lstlisting}

\begin{lstlisting}[caption={Wyznaczanie wspólczynników aproksymacji Czebyszewa}]
N_quad = 500
j_idx  = np.arange(1, N_quad + 1)
theta_j = (2 * j_idx - 1) * np.pi / (2 * N_quad)
t_j = np.cos(theta_j)      # wezly Gaussa-Czebyszewa
x_j = t_j + 1              # przedzial [0, 2]
f_j = np.sqrt(x_j)

deg = 2
c = np.zeros(deg + 1)
c[0] = np.sum(f_j) / N_quad
for k in range(1, deg + 1):
    c[k] = 2.0 / N_quad * np.sum(f_j * np.cos(k * theta_j))
\end{lstlisting}

% --- 2.3 Wyniki ---
\subsection{Wyniki}

\noindent
Wyznaczone współczynniki (w~bazie Czebyszewa, na~$[-1, 1]$):
\[
  c_0 \approx 0{,}9003, \quad c_1 \approx 0{,}6002, \quad
  c_2 \approx -0{,}1200.
\]

\noindent
Wielomian w~postaci standardowej:
\[
  p(x) \approx 0{,}1801 + 1{,}0804 \cdot x - 0{,}2401 \cdot x^2.
\]

\begin{figure}[H]
  \centering
  \includegraphics[width=0.95\textwidth]{wykres4.png}
  \caption{Aproksymacja $f(x) = \sqrt{x}$ wielomianem 2. stopnia
           (baza Czebyszewa) i~wykres błędu bezwzględnego.}
  \label{fig:cheb_approx}
\end{figure}

\noindent
Błędy aproksymacji:
\begin{itemize}
  \item Maksymalny błąd bezwzględny: $\|f - p\|_\infty \approx 1{,}80 \times 10^{-1}$
  \item Błąd $L_2 \approx 5{,}21 \times 10^{-2}$
\end{itemize}

\subsubsection{Porównanie z~klasyczną metodą MNK}

\begin{figure}[H]
  \centering
  \includegraphics[width=0.95\textwidth]{wykres5.png}
  \caption{Porównanie aproksymacji Czebyszewa z~klasycznym wielomianem
           MNK stopnia~2 oraz wykresy błędów obu metod.}
  \label{fig:mnk_vs_cheb}
\end{figure}

% --- 2.4 Wnioski ---
\subsection{Wnioski}

\begin{enumerate}
  \item Baza Czebyszewa daje \emph{aproksymację quasi-jednostajną},
        ponieważ minimalizuje maksymalne odchylenie od funkcji
        globalnie na~całym przedziale.

  \item Standardowa aproksymacja MNK minimalizuje sumę kwadratów
        odchyleń (norma~$L_2$), co może prowadzić do większych błędów
        w~skrajnych punktach przedziału w~porównaniu z~metodą
        Czebyszewa.

  \item Metoda Czebyszewa jest numerycznie stabilna ze~względu na
        ortogonalność bazy, a współczynniki wyznacza się niezależnie,
        bez rozwiązywania układu równań.

  \item Aproksymacja wielomianami Czebyszewa jest \emph{tania
        obliczeniowo}, ponieważ zamiast rozwiązywać układ $m \times m$,
        obliczamy $m$ niezależnych całek.

  \item Dla $m = 2$ błąd $L_2 \approx 5{,}21 \times 10^{-2}$, co
        jest dobrym wynikiem dla wielomianu tak niskiego stopnia
        na~przedziale $[0, 2]$.
\end{enumerate}

% ============================================================
\section{Podsumowanie}
% ============================================================

\begin{enumerate}
  \item \textbf{Aproksymacja a interpolacja.} Aproksymacja nie musi
        przechodzić przez węzły, dzięki czemu jest odporna na szum
        w~danych i unika problemów z przeuczeniem.

  \item \textbf{Wybór stopnia wielomianu.} Zbyt niski stopień
        (underfitting) oznacza, że model nie opisuje trendu. Zbyt
        wysoki (overfitting) oznacza, że model uczy się szumu i
        źle generalizuje. Kryterium AICc pozwala automatycznie
        wybrać kompromis bez potrzeby posiadania zestawu testowego.

  \item \textbf{Baza Czebyszewa} jest lepszym wyborem niż baza
        potęgowa dla aproksymacji ciągłej, ponieważ jest stabilna
        numerycznie, daje aproksymację quasi-jednostajną i może
        zastąpić kosztowną aproksymację minimax.

  \item \textbf{Normalizacja dziedziny} (np.~liniowe przekształcenie
        do $[-1, 1]$) jest niezbędna, by uniknąć złego uwarunkowania
        macierzy przy użyciu bazy potęgowej lub Czebyszewa.
\end{enumerate}

\end{document}
